\chapter{Results and Comparative Evaluation}

\section{Chapter Introduction}
This chapter presents the reported evaluation evidence and explains how the remaining comparative analysis must be completed. It separates measured or reported findings from completion instructions so that the thesis does not claim results that have not yet been produced. The chapter also adds explicit guidance for HR@k and NDCG@k evaluation, because these ranking metrics are necessary for a defensible retrieval-comparison section.

\section{Reported Experimental Results}
This chapter presents the reported evaluation findings using two sources of evidence: the internships.click extraction output and a multi-site job-search simulation. It summarizes the performance of the extraction and classification modules, discusses retrieval quality, and analyzes the system's operational efficiency. Because the available evidence is partly narrative and not accompanied by complete benchmark tables, the findings are interpreted as reported experimental evidence rather than as a fully reproducible external-platform benchmark. The chapter is organized as follows:


\begin{itemize}[leftmargin=1.4em,itemsep=0.25em,topsep=0.25em]

\item The extraction-results subsection reports on Layer A parsing of structured fields from the internships.click dataset.

\item The classification-results subsection evaluates semantic mapping from raw job titles to canonical roles.

\item The retrieval-results subsection analyzes aggregation performance across multiple job platforms.

\item The efficiency-and-cost subsection examines the operational implications of the tiered architecture.

\end{itemize}
\subsection{Extraction Results}
The first stage of evaluation used the internships.click output dataset, where job fields such as company, location, and salary were often embedded within a single raw text block, for example ``Matrox Graphics Inc Dorval \$95K - \$124K.'' The reported results indicate that the rule-based Layer A reached 98\% precision for compensation extraction, identifying patterns such as ``\$95K - \$124K'' and ``\$23.00/hour.'' Recall for explicit location fields was lower, reported at 76\%, when geographic markers were ambiguous within the title string. These findings suggest that deterministic patterns are stable for highly regular fields [44], while semantic fallback is still required for higher coverage.

\subsection{Classification Results}
Layer B was evaluated by its ability to normalize raw titles from the internships.click data, such as ``Spring 2026 Electronics Software Development Intern (MS),'' into canonical occupations. The reported Top-3 accuracy was 91\%. The qualitative analysis indicates that, even when the exact label was missed, the bi-encoder architecture [42] often mapped roles to the appropriate professional family, such as mapping ``General Software Developer Intern- Humanoid Robot'' to ``Software Engineer.'' This result indicates an improvement over keyword matching for titles with diverse employer-specific phrasing.

\subsection{Retrieval Results}
Retrieval quality was assessed using the multi-site job simulation, which aggregated results across four platforms: We Work Remotely, Palantir (Lever), Saronic (Lever), and Box (Greenhouse). The reported simulation included 416 unique listings across five search categories: Software, AI, Data, Web, and Developer. The integration of dense embeddings allowed the system to retrieve relevant results even when keyword overlap was limited, such as retrieving ``Forward Deployed AI Engineer'' roles for ``Software'' queries. The reported comparison indicates a +22\% gain in P@5 relative to lexical-only baselines, although the final thesis still requires the complete query set and relevance judgments for full reproducibility.

\subsection{Student Procedure for HR@k and NDCG@k Analysis}
The final version of this chapter must include a reproducible retrieval-ranking evaluation. Students should complete this part as a measured experiment, not as a narrative comparison. The following procedure should be followed before inserting numerical values in the thesis:

\begin{enumerate}[leftmargin=1.7em,itemsep=0.25em,topsep=0.25em]

\item Freeze the query set before running the comparison. The same queries must be used for the proposed system, the keyword-only baseline, and any external platform included in the comparative analysis.

\item For each query and each system, export the top 10 results with rank, title, company, source platform, URL, retrieval score if available, and the date on which the result was captured.

\item Judge each query-result pair using the same relevance scale. A practical three-level scale is 0 for not relevant, 1 for partially relevant, and 2 for highly relevant. The criteria must be written before scoring begins.

\item Convert the judged lists into ranking metrics. HR@5 and HR@10 should record whether at least one relevant result appears in the first 5 or 10 results. NDCG@10 should use the graded relevance labels, so that highly relevant results placed near the top receive more credit than weakly relevant or lower-ranked results.

\item Report the mean value of each metric across the complete query set. Do not compare systems on different query sets, different capture dates, or incomplete result lists.

\item Insert screenshots only as qualitative evidence. Screenshots should support the numerical table; they should not replace HR@k, NDCG@k, MRR, MAP, or Precision@k.

\end{enumerate}

The table below is a completion template. It must remain empty or marked as not measured until the students compute the values from the frozen query set and relevance judgments.

\begingroup
\small
\begin{longtable}{p{0.18\textwidth}p{0.18\textwidth}p{0.08\textwidth}p{0.08\textwidth}p{0.08\textwidth}p{0.08\textwidth}p{0.09\textwidth}}
\caption{Template for retrieval-ranking metrics to be completed with measured evidence.}\\
\toprule
\textbf{System} & \textbf{Query group} & \textbf{HR@5} & \textbf{HR@10} & \textbf{MRR} & \textbf{MAP} & \textbf{NDCG}\newline\textbf{@10}\\
\midrule
Proposed hybrid system & To be measured & -- & -- & -- & -- & --\\
Keyword-only baseline & To be measured & -- & -- & -- & -- & --\\
External platform comparison & To be measured & -- & -- & -- & -- & --\\
\bottomrule
\end{longtable}
\endgroup

\subsection{Efficiency and Cost}
The computational metrics from the multi-site simulation provide evidence for the feasibility of the tiered pipeline. The reported configuration used Layer B for 95.4\% of the 416 simulated listings and reserved Layer C for the most ambiguous cases. Under this reported configuration, LLM API cost was reduced by 92\% relative to an LLM-only approach. This evidence supports the Green AI [46] objective of reducing unnecessary model calls while maintaining usable extraction and classification quality.



\section{Comparison with Other Platforms and Baselines}
This section is retained as an organizational scaffold for the final thesis. The original thesis discusses the need to compare keyword-based search, LLM-dependent processing, general recruitment platforms, and university career portals, but it does not provide an independently measured external-platform benchmark. Therefore, the comparison below does not introduce new empirical findings; it records what evidence must be inserted when the student completes the evaluation.

\begin{longtable}{p{0.24\textwidth}p{0.34\textwidth}p{0.32\textwidth}}
\caption{Comparison scaffold to be completed with measured evidence.}\\
\toprule
\textbf{Comparison target} & \textbf{Purpose in the thesis} & \textbf{Evidence still required}\\
\midrule
Keyword-only internship search & Test whether the proposed hybrid retrieval approach improves over exact lexical matching. & Same query set, keyword-only results, judged relevance, and HR@k, Precision@k, MRR, MAP, and NDCG@k tables.\\
LLM-only enrichment pipeline & Test the cost and latency trade-off against a system that delegates extraction and normalization to an external LLM. & Cost per 1,000 listings, latency distribution, failure examples, and a clear prompt/protocol.\\
General recruitment platforms & Contextualize the proposed platform against mature public internship/job-search systems. & Side-by-side screenshots for identical queries and manually judged relevance.\\
University career portals & Compare against the institutional use case targeted by the thesis. & Partner JSON upload examples, before/after normalization outputs, and search screenshots.\\
\bottomrule
\end{longtable}

\section{Interpretation of the Results}
The interpretation in this section is limited to the results already reported in the original thesis. The current evidence supports an architectural argument for a staged semantic pipeline, but the final comparison requires measured tables, screenshots, and reproducible benchmark artifacts before it can be treated as a complete empirical comparison.

\section{Chapter Conclusion}
This chapter has reported the evidence currently available and has defined the missing retrieval-comparison evidence that must be completed before final defense. The added HR@k and NDCG@k procedure ensures that future result values will be based on a fixed query set, explicit relevance judgments, and reproducible ranked lists rather than unsupported comparative claims.
